\documentclass[11pt]{article}
\usepackage[margin=1in]{geometry}
\usepackage{amsmath,amssymb,amsthm,hyperref}
\newtheorem{proposition}{Proposition}
\newcommand{\E}{\mathbb E}
\newcommand{\Prob}{\mathbb P}
\title{Square-root arity and approximation across a split of the hash seeds}
\author{Project research note}
\date{18 November 2026}
\begin{document}
\maketitle

Let $b=\lceil\log_2p\rceil$ and $N=\#E(\mathbb F_p)$, supplied as $hr$
with $r$ an odd prime and $\gcd(h,r)=1$. The best complete bound in the
project remains the conditional
\[
 \log_2\E T\le (26.2933775156+o(1))\,b/(\log_2 b)^{1/4}.
\]
It uses the packet assumption A.1 and the arithmetic, contraction and
relation-reduction results specified in \texttt{../short\_hash\_count/count.tex}.
The original paper gives the same asymptotic scale on the common
prime-order class without a numerical leading constant
(\href{https://github.com/bochum-crypto/subexponential-ecdl/blob/main/subexponential-ecdl.pdf}
{Akshima et al., version 0.2, Theorems 1, 2 and 4}).

\section{The outer cost at square-root arity}
Put $m=\lceil\sqrt b\rceil$, $M=2^b$. For independently hashed slots let
$w_i$ be their dyadic widths, $\rho_i=w_i/M$, $\nu_i=N\rho_i$, and
$\lambda=N^{-1}\prod_i\nu_i$. For a fixed target let $H$ be the
ordinary integer tuple count and $X$ the count obeying the chord guards.
Take the balanced isolation widths from \texttt{mixed.tex}, so
$1/3<\lambda\le2/3$. The sum of the integer width exponents is
$\lfloor\log_2((2/3)M^m/N^{m-1})\rfloor$, and they differ by at most one.
Consequently $\log_2\nu_i=(\log_2 N)/m+O(1)$ uniformly.
Hasse's bound gives $\log_2 N=b+O(1)$.

For $2T\ne O$, the finite factorial-moment inequality of \texttt{mixed.tex} gives
\[
 \E[H(H-1)]\le\lambda^2\prod_i(1+2/\nu_i)
       +\lambda |E(\mathbb F_p)[2]|(2^m-2)/N,\qquad
 \E X\ge(1-3m/N)\lambda .
\]
Here $m/\min_i\nu_i=o(1)$ and $2^m/N=o(1)$ also at $m=\lceil\sqrt b\rceil$.
Thus $\Prob(X=H=1)\ge\lambda(1-\lambda)-o(1)\ge2/9-o(1)$.
All masked nonidentity targets for cofactor projection are in the
odd-prime-order subgroup, where $2T\ne O$. The probability is over
fresh independent seeds for each prescribed target.

\begin{proposition}
Suppose ordinary integer-count parity on every such independently
seeded input and every legal dyadic restriction has an exact evaluator
with expected bit cost $2^{\widetilde O(\sqrt b)}$ per call,
including preprocessing and uniformly conditional on the input and
the preceding transcript. The project outer reduction then gives
$\log_2\E T=\widetilde O(\sqrt b)$.
\end{proposition}
\begin{proof}
On odd parity, the restriction argument in \texttt{mixed.tex}
uses at most $1+\sum_i\log_2 w_i=O(b+m)$ calls including the root,
and at most $2^m$ terminal sign choices. Deterministic selection and a
guard check give the same relation-selection law in training and online
use. An isolated guarded witness always succeeds. The enumeration
budget is $B=2\sum_i w_i$, with $\log_2 B=b/m+O(\log m+1)$.
The pilot and sparse-solve bounds charge polynomial factors in $b,m$
times $B C_m+mB^2$, where $C_m$ includes the full selection procedure.
Their proofs use a dimension bound $d\le B$, $O(m)$ row weight, and a
constant average success probability. The small-subgroup fallback costs
$B\operatorname{poly}(b,m)$. Each term fits the proposed bound. This
extends the outer estimates individually without extrapolating the
contraction theorem.
\end{proof}
At polylogarithmic arity the current estimate is
$\log_2 C_m\le(A+o(1))mb/\sqrt{\log_2b}$, $A=18K$ in the short-hash
note. Reoptimizing the two full-cost terms gives
$m\sim(\log_2 b)^{1/4}/\sqrt A$. Even a uniform extension of this
estimate to square-root arity would give a term of order
$b^{3/2}/\sqrt{\log_2 b}$. For a smaller strict gain it would suffice,
for example, to save a factor $g(b)=\log_2\log_2 b$ in the relation
exponent uniformly at $m\sim(\log_2 b)^{1/4}\sqrt{g(b)/A}$. The two
full-cost terms then give $(2\sqrt A+o(1))b/((\log_2 b)^{1/4}\sqrt{g(b)})$.
Neither evaluator bound follows from the results below.

\section{Optimal separated approximation at a fixed target}
Here $h=1$, $r=N$, $T\ne O$, and each slot has an independent uniform
odd multiplier and uniform offset modulo $M$. A predetermined dyadic
subinterval of a slot has this same law by translation of its offset.
Split the slots into any two nonempty parts. A rank-$q$ real
approximation has the form $\sum_{j=1}^q a_j(\text{left seeds})
b_j(\text{right seeds})$. Its squared error is averaged over the full
fresh seed law. Write $\|Y\|_2^2=\E Y^2$.

\begin{proposition}
Let $r\to\infty$, $m/r+\sum_i\rho_i=o(1)$, and $n=(r-1)/2$.
For $P=\prod_i|F_i|/r$, $R=H-P$, and $q/n\to x\in[0,1]$,
\[
 \inf_{\operatorname{rank} Z\le q}\frac{\|R-Z\|_2^2}{\lambda}
 \ \longrightarrow\ 1-x-\frac{\sin(\pi x)}{\pi}.
\]
At sparse balanced widths with $\lambda=o(1)$, $m/\min_i\nu_i=o(1)$
and $2^m/r=o(1)$, the same limit holds with $R$ replaced by either
$V=H\bmod 2$ or $V=1_{\{X>0\}}$ and denominator $\E V^2$.
For example, at $m=\lceil\sqrt b\rceil$ choose balanced dyadic widths
whose product belongs to $(\sqrt r/2,\sqrt r]$. All the sparse-width
conditions hold.
\end{proposition}
\begin{proof}
Use the independent-seed Fourier representation in \texttt{mixed.tex}.
In probability-weighted row and column spaces the residual is
$\sqrt{\lambda/r}\,U D V_0^*$, where $U,V_0$ have operator norm at
most one and $D$ has diagonal entries $2\cos(2\pi jT/r)$ for one
representative of each nontrivial opposite pair. Here $T$ in the
cosine denotes its cyclic coordinate, used only in the proof.
The same Gram-matrix identity puts the squared Frobenius norm of
$U D V_0^*$ between $\|D\|_F^2-\Delta$ and $\|D\|_F^2$, with
$\Delta=4\sum_i(1+n\rho_i)=o(r)$.

Let $s_1\ge\cdots\ge s_n$ be the sorted absolute diagonal entries.
The $j$th singular value of $U D V_0^*$ is at most $s_j$.
After subtracting its top $q$ squared singular values the squared
norm remaining is at least $\sum_{j>q}s_j^2-\Delta$.
Keeping the $q$ largest entries of $D$ inside the product gives a
rank-$q$ approximation of squared error at most $\sum_{j>q}s_j^2$
in the same normalization. Thus the bounds differ by $o(r)$,
uniformly in $q$. Multiplication by the target permutes opposite
pairs. The squared diagonal entries sample $4\cos^2\theta$ at
angles asymptotically uniform on $[0,\pi]$. Integrating the fraction
$x$ nearest the endpoints gives captured fraction
$x+\sin(\pi x)/\pi$. The total squared norm is $r-2$.

The factorial bound implies $\E H(H-1)=o(\lambda)$ at the sparse
widths. Also $\E(H-X)\le3m\lambda/r$. Pointwise, for either bit,
\[
 (H-V)^2\le2H(H-1)+1_{\{H=1,X=0\}} .
\]
The isolation bound and $\E H\le\lambda$ give
$\E V^2=(1+o(1))\lambda$. The two-abscissa hash bound gives
$\operatorname{Var}|F_i|\le2\E|F_i|$ on a prime-order curve.
Independence yields
$\E P^2\le\lambda^2\prod_i(1+2/((r-1)\rho_i))=o(\lambda)$.
Thus $\|V-R\|_2=o(\sqrt{\lambda})$. Distance to any fixed rank class
is 1-Lipschitz in this norm, proving the assertion also at $x=0,1$.
\end{proof}
Half of the opposite pairs leave normalized squared error $1/2-1/\pi$.
Any $o(r)$ separated terms leave asymptotically all the squared error
of the zero predictor in the sparse decision problem. An attaining
mathematical expansion can be expensive to evaluate. These limits
concern the real approximation itself under fresh seeds. Thresholding
a score, conditioning on a selected dyadic path, and computing jointly
across seeds have different hypotheses. Rare positives dominate this
relative error.

\section{Two finite bounds}
\paragraph{Fixed multipliers.}
The preceding spectral formula averages over the multipliers. A weaker
bound applies for \emph{every} fixed vector of odd multipliers. Let
$W=\prod_iw_i$. Write the guarded count at $T\ne O$ as a matrix $A$ over
the offsets in a split with $l,u\ge1$ slots, $l+u=m$, on a prime-order
curve. Each generic tuple occurs for exactly $W$ offset vectors.
Consequently $\|A\|_F^2\ge\sum A\ge(1-3m/r)r^{m-1}W$.
At fixed left offsets there are at most $\prod_{i\in L}(2w_i)$
left tuples. Each has at most $r^{u-1}$ right completions per tuple
when the points are unrestricted, each appearing at most
$\prod_{i\in R}w_i$ times across the right offsets.
The largest row sum is at most $2^lWr^{u-1}$ and the largest column
sum at most $2^uWr^{l-1}$. Their product bounds $\|A\|_{\rm op}^2$.
A rank-$q$ approximation of relative Frobenius error $\varepsilon<1$
therefore requires
\[
 q\ \ge\ (1-\varepsilon^2)(1-3m/r)\frac{r}{2^mW}.
\]
The same proof applies to the ordinary count. At $W\asymp\sqrt r$,
$m=\lceil\sqrt b\rceil$, this is $2^{b/2-O(\sqrt b)}$ terms.
The claim is about relative mean-square error at fixed multipliers
and fresh offsets, including predetermined subwindows.

\paragraph{A smaller sufficient width for constant coverage.}
Return to general $N$ and a fixed target with $2T\ne O$. Let $F=N\rho$
for equal widths, $d_2=|E(\mathbb F_p)[2]|\le4$, $\gamma=1-3m/N$,
and $\mu=\E X$. The original outer coverage lemma prices all equal-abscissa
tuples together. Using the mixed-note sign analysis separately yields
\[
 \frac{\E X^2}{\mu^2}\le
 \frac{(1+2/F)^m-(2/F)^m+N/F^m}{\gamma}
 +\frac{d_2(2^m-2)}{\gamma^2 F^m}.
\]
For a proper set of $k$ matching abscissas, use the actual number of
generic first tuples, at most $2^kN^{m-k-1}$ second tuples per first
tuple, and the two-abscissa hash bound. Summing and dividing by
$\mu^2$ gives the first two binomial terms. For a complete match,
the identical tuple contributes $\mu$ before normalization. Every
other tuple flips a nonempty proper set of signs since $2T\ne O$.
The sum on the flipped coordinates is in $E[2]$, giving at most
$d_2N^{m-2}$ choices per flip set and joint acceptance $\rho^m$.
Substitution of $\mu\ge\gamma F^m/N$ proves the formula.

For $N\ge30m$, $m\ge2$ and $F\ge\max(N^{1/m},20m)$, the right side
is at most $(e^{1/10}+1)/0.9+4/(20^2\cdot0.9^2)<2.5$.
Thus $\Prob(X>0)>2/5$. Relative to the previous sufficient threshold
$\max(2N^{1/m},20m)$ this can save at most one dyadic width bit.
It does not change the leading coefficient or the asymptotic scale.

\paragraph{Finite checks.}
Run \texttt{python -B research/mixed\_coefficient/check\_frontier.py}.
It checks singular-value tail bounds on two small prime-order curves,
both decision-bit comparisons, and the limiting cosine integral.
\end{document}
